
\section*{4. Passage to the quotient by a finite flat groupoid
  (proof of a special case)}
\phantomsection\label{sec:sga3-V-4}
\addcontentsline{toc}{section}{4. Passage to the quotient by a finite flat groupoid}

\medskip
\phantomsection\label{thm:sga3-V-4.1}
\noindent\textbf{Theorem 4.1.}---
\textit{Consider a \((\mathbf{Sch}/S)\)-groupoid}

\begin{equation*}
  \resizebox{0.40\textwidth}{!}{%
    \SGAThreeNativeInput{figures/exp5/p261_theorem41_groupoid.png}}
  \phantomsection\label{fig:sga3-exp5-theorem41-groupoid}
\end{equation*}

\textit{and suppose that the following conditions hold:}%
\begin{enumerate}
\item[\textit{a)}]\phantomsection\label{thm:sga3-V-4.1-a}%
  \textit{\(d_1\) is finite locally free;}
\item[\textit{b)}]\phantomsection\label{thm:sga3-V-4.1-b}%
  \textit{for every \(x\in X_0\), the set
  \(d_0d_1^{-1}(x)\) is contained in an affine open subset of
  \(X_0\).}%
\end{enumerate}
\textit{Then:}
\begin{enumerate}
\item[(i)]\phantomsection\label{thm:sga3-V-4.1-i}%
  \textit{there exists a cokernel \((Y,p)\) of \((d_0,d_1)\) in
  \((\mathbf{Sch}/S)\). Moreover, such a pair \((Y,p)\) is a
  cokernel of \((d_0,d_1)\) in the category of all ringed spaces;}
\item[\SGAInlineRefTarget{thm:sga3-V-4.1:item:ii}(ii)]\phantomsection\label{thm:sga3-V-4.1-ii}%
  \textit{\(p\) is integral and open, and \(Y\) is affine when
  \(X_0\) is affine;}%
  \footnote{\label{note:sga3-V-18}Editor's note: the assertion that
  \(p\) is open has been added by repeating the analogous argument
  given in \hyperref[lem:sga3-V-6.1]{6.1}.}
\item[(iii)]\phantomsection\label{thm:sga3-V-4.1-iii}%
  \textit{the morphism
  \[
    X_1\longrightarrow X_0\times_YX_0
  \]
  with components \(d_0\) and \(d_1\) is surjective;}
\item[(iv)]\phantomsection\label{thm:sga3-V-4.1-iv}%
  \textit{if \((d_0,d_1)\) is an equivalence pair, then
  \(X_1\to X_0\times_YX_0\) is an isomorphism}%
  \textit{and \(p:X_0\to Y\) is finite locally free.}%
  \footnote{\label{note:sga3-V-20}Editor's note: the consequences that
  follow have been made explicit; see \(\hyperref[bib:sga3-V-Ray67a]{[{\rm Ray67a}]}\),
  Theorem~1(iii), and \hyperref[proof:sga3-V-5-b-general]{the proof given below} at the end of \hyperref[sec:sga3-V-5]{Section~5}.}
  \textit{Moreover, \((Y,p)\) is a cokernel of \((d_0,d_1)\) in the
  category of sheaves for the fppf topology and, for every base change
  \(Y'\to Y\), the scheme \(Y'\) is the cokernel of the groupoid
  \(X_\ast\times_YY'\) obtained from \(X_\ast\) by the base change
  \(X_0\times_YY'\to X_0\).}
\end{enumerate}
\textit{In particular, for every base change \(S'\to S\), the scheme
\(Y'=Y\times_SS'\) is the cokernel of the \(S'\)-groupoid
\(X'_\ast=X_\ast\times_SS'\). Thus, in this case, formation of the
quotient commutes with base change.}

It follows immediately from \hyperref[thm:sga3-V-4.1-i]{(i)} that the underlying topological space
of \(Y\) is the quotient of the underlying topological space of
\(X_0\) by the equivalence relation defined by the
\((\mathbf{Sch}/S)\)-groupoid \(X_\ast\).

We shall first prove \hyperref[thm:sga3-V-4.1]{the theorem} when \(X_0\) is affine and \(d_1\) is
locally free of constant rank \(n\). We shall then see how to reduce
the general situation to this special case.
